The list of inputs is Fact \claimref{dark-matter.imported-inputs-standard-model-chain}, whose own lemmas (pair count, Segre) are statements about $V_{\C}$ and about tensor products of complex spaces. Invariance item by item.

Complexification: $\R^{2,4} \otimes \C = \R^{4,2} \otimes \C$ as quadratic spaces, there being no signature over $\C$. Even subalgebra: $e_{i} \mapsto e_{i}e_{n}$ ($i < n$) identifies $\Cl^{0}(p,q)$ with $\Cl(p, q-1)$ when $e_{n}^{2} = -1$; the script verifies the relations exactly for $(2,4)$ and $(4,2)$, and Theorem \claimref{dark-matter.type-residue-what-each-tower-can}'s computation gives $\Cl(2,3)$ and $\Cl(4,1)$ the same type, centre and module dimension, both $\Mat(4,\C)$. The spin group (even versors with $R\tilde{R} = 1$) and the invariant form $\bar{\psi}\varphi$ on $\Cl^{0}f$ are defined by the even subalgebra and the ambient reversion, which the reversal leaves unchanged, since $(p,q) \mapsto (q,p)$ changes the signs of vector squares and not the order of products. Under $e_{i} \mapsto e_{i}e_{n}$ the ambient reversion becomes Clifford conjugation of $\Cl(2,3)$, the antilinear involution a hermitian form needs, and the form so built has signature $(2,2)$: the spin group is $SU(2,2)$ in both cases, with the same representation on the conformal spinor space $\CS$, the minimal left ideal of the even algebra (likewise $\Cl^{0}(1,3) \cong \Cl(1,2)$ and $\Cl^{0}(3,1) \cong \Cl(3,0)$ are both $\Mat(2,\C)$, with the same Lorentz spin group and Weyl spinors).

Mass formula: $m^{2} = 2|\br{\pi^{1}\pi^{2}}|^{2}$ is the squared norm of $\pi^{1} \wedge \pi^{2}$ for Weyl spinors of the common $SL(2,\C)$. Segre and hypercharge: statements about $\C^{2} \otimes \C^{3}$ and mode counts. Hodge star: on two-forms in four dimensions
\[
(\star F)_{ab} \;=\; \tfrac{1}{2}\sqrt{|g|}\,\epsilon_{abcd}\,g^{ce}g^{df}F_{ef},
\]
even in $g$, so $F \wedge \star F$ and $F \wedge F$ are reversal-invariant.

(i) follows by running the paper's chain on the reversed data; (ii) is the observation that the real module of $\Cl(p,q)$ appears nowhere in the list. (iii): identical content gives identical $\beta$-functions in any scheme, to all orders; with $\alpha_{2}(M) = \alpha_{6}(M)$ the running solutions coincide, so $\Lambda_{2} = \Lambda_{6}$, and the composite spectra, being those of isomorphic theories, coincide, $c_{2} = c_{6}$; then, under P4,
\[
f_{2} \;=\; \frac{\epsilon_{2}}{\epsilon}\cdot\frac{n_{2}}{n_{6}}\cdot\frac{m_{2}}{m_{6}} \;=\; 1 \cdot 1 \cdot 1 .
\] (iv) is the audit clause of Premise \claimref{dark-matter.p2-indifference} read in the direction it permits.

The second route of (v) rests on \claimref{dark-matter.p1-single-frozen-point}, whose orientation clause orients both towers by the one Witt flag of the descent.
